
\documentclass{article}
\usepackage{colortbl}
\usepackage{makecell}
\usepackage{multirow}
\usepackage{supertabular}

\begin{document}

\newcounter{utterance}

\centering \large Interaction Transcript for game `codenames', experiment `opponent\_none', episode 2 with glm{-}5.2.
\vspace{24pt}

{ \footnotesize  \setcounter{utterance}{1}
\setlength{\tabcolsep}{0pt}
\begin{supertabular}{c@{$\;$}|p{.15\linewidth}@{}p{.15\linewidth}p{.15\linewidth}p{.15\linewidth}p{.15\linewidth}p{.15\linewidth}}
    \# & $\;$A & \multicolumn{4}{c}{Game Master} & $\;\:$B\\
    \hline

    \theutterance \stepcounter{utterance}  
    & & & \multicolumn{2}{p{0.3\linewidth}}{
        \cellcolor[rgb]{0.9,0.9,0.9}{
            \makecell[{{p{\linewidth}}}]{
                \texttt{\tiny{[GM$|$GM]}}
                {'hidden': {'team': ['Hai', 'Stamm', 'Krawatte', 'Rücken', 'Bildschirm', 'Zelle', 'Schwingen', 'Wirbelsäule', 'Hufeisen'], 'innocent': ['Schalter', 'Tonhöhe', 'China', 'Lehrer', 'Tokio', 'Zahn', 'Drachen'], 'opponent': ['Hafen', 'Socke', 'Mitarbeiter', 'Mütze', 'Haken', 'Truthahn', 'Umhang', 'Skala'], 'assassin': ['Schule']}, 'revealed': {'team': {'team': [], 'innocent': [], 'opponent': [], 'assassin': []}, 'opponent': {'team': [], 'innocent': [], 'opponent': [], 'assassin': []}}}
            }
        }
    }
    & & \\ \\

    \theutterance \stepcounter{utterance}  
    & & \multicolumn{4}{p{0.6\linewidth}}{
        \cellcolor[rgb]{0.9,0.9,0.9}{
            \makecell[{{p{\linewidth}}}]{
                \texttt{\tiny{[P1$\langle$GM]}}
                \texttt{Gib einen einzigen einwortigen Hinweis, der zu einigen der Wörter deines Teams aus der Liste unten passt. Du kannst wählen, zu wie vielen und zu welchen Wörtern sich dein Hinweis bezieht.} \\
\texttt{Das Ziel ist, dass dein*e Partner*in anhand der Hinweise, die du gibst, alle Wörter deines Teams errät, bevor das andere Team alle seine aufgedeckt hat. Wenn dein*e Partner*in ein "Gegner"{-} oder "Ablenker"{-}Wort errät, endet sein*ihr Ratezug und das gegnerische Team ist am Zug. Wenn dein*e Partner*in ein "Killer"{-}Wort errät, endet das Spiel und dein Team verliert sofort.} \\
\texttt{Je mehr Wörter sich dein Hinweis bezieht, desto schneller kannst du das Spiel gewinnen, aber desto schwieriger können die Assoziationen für dein*e Partner*in zu erraten sein – wähle dein Hinweiswort also mit Bedacht. Das Hinweiswort muss semantisch mit den Zielwörtern verwandt sein; es darf keines der Wörter in den Listen sein oder Teile davon enthalten.} \\
\\ 
\texttt{Gib immer deinen einwortigen Hinweis und deine kommagetrennte Liste der zugehörigen Zielwörter im folgenden Format an und halte deine Antworten so kurz wie möglich. Füge niemals anderen Text hinzu, als in dieser Form erforderlich ist:} \\
\\ 
\texttt{HINWEIS: <WORT>} \\
\texttt{ZIELE: <ZIELE>} \\
\\ 
\texttt{Die Wörter deines Teams sind: Hai, Stamm, Krawatte, Rücken, Bildschirm, Zelle, Schwingen, Wirbelsäule, Hufeisen.} \\
\texttt{Die Wörter des gegnerischen Teams sind: Hafen, Socke, Mitarbeiter, Mütze, Haken, Truthahn, Umhang, Skala.} \\
\texttt{Ablenker{-}Wörter sind: Schalter, Tonhöhe, China, Lehrer, Tokio, Zahn, Drachen.} \\
\texttt{Die Killer{-}Wörter sind: Schule.} \\
            }
        }
    }
    & \\ \\

    \theutterance \stepcounter{utterance}  
    & \multicolumn{4}{p{0.6\linewidth}}{
        \cellcolor[rgb]{0.9,0.9,0.9}{
            \makecell[{{p{\linewidth}}}]{
                \texttt{\tiny{[P1$\rangle$GM]}}
                \texttt{HINWEIS: Anatomie} \\
\texttt{ZIELE: Rücken, Wirbelsäule, Zelle, Stamm} \\
            }
        }
    }
    & & \\ \\

    \theutterance \stepcounter{utterance}  
    & & & \multicolumn{2}{p{0.3\linewidth}}{
        \cellcolor[rgb]{0.9,0.9,0.9}{
            \makecell[{{p{\linewidth}}}]{
                \texttt{\tiny{[GM$|$GM]}}
                \texttt{anatomie} \\
            }
        }
    }
    & & \\ \\

    \theutterance \stepcounter{utterance}  
    & & & \multicolumn{2}{p{0.3\linewidth}}{
        \cellcolor[rgb]{0.9,0.9,0.9}{
            \makecell[{{p{\linewidth}}}]{
                \texttt{\tiny{[GM$|$GM]}}
                ['rücken', 'wirbelsäule', 'zelle', 'stamm']
            }
        }
    }
    & & \\ \\

    \theutterance \stepcounter{utterance}  
    & & & \multicolumn{2}{p{0.3\linewidth}}{
        \cellcolor[rgb]{0.9,0.9,0.9}{
            \makecell[{{p{\linewidth}}}]{
                \texttt{\tiny{[GM$|$GM]}}
                {'word': 'rücken', 'assignment': 'team'}
            }
        }
    }
    & & \\ \\

    \theutterance \stepcounter{utterance}  
    & & & \multicolumn{2}{p{0.3\linewidth}}{
        \cellcolor[rgb]{0.9,0.9,0.9}{
            \makecell[{{p{\linewidth}}}]{
                \texttt{\tiny{[GM$|$GM]}}
                {'word': 'wirbelsäule', 'assignment': 'team'}
            }
        }
    }
    & & \\ \\

    \theutterance \stepcounter{utterance}  
    & & & \multicolumn{2}{p{0.3\linewidth}}{
        \cellcolor[rgb]{0.9,0.9,0.9}{
            \makecell[{{p{\linewidth}}}]{
                \texttt{\tiny{[GM$|$GM]}}
                {'word': 'zelle', 'assignment': 'team'}
            }
        }
    }
    & & \\ \\

    \theutterance \stepcounter{utterance}  
    & & & \multicolumn{2}{p{0.3\linewidth}}{
        \cellcolor[rgb]{0.9,0.9,0.9}{
            \makecell[{{p{\linewidth}}}]{
                \texttt{\tiny{[GM$|$GM]}}
                {'word': 'stamm', 'assignment': 'team'}
            }
        }
    }
    & & \\ \\

    \theutterance \stepcounter{utterance}  
    & & \multicolumn{4}{p{0.6\linewidth}}{
        \cellcolor[rgb]{0.9,0.9,0.9}{
            \makecell[{{p{\linewidth}}}]{
                \texttt{\tiny{[GM$\rangle$P2]}}
                \texttt{Gib eine durch Kommas getrennte Liste mit bis zu 4 Wörtern aus der folgenden Liste an, die sich am besten auf das Wort "anatomie" beziehen oder am engsten damit verbunden sind. Beginne deine Liste der Vermutung(en) immer mit TIPP:  und füge keine anderen Texte in deine Antwort ein.} \\
\\ 
\texttt{Schule, Tokio, Krawatte, Mitarbeiter, Hai, Haken, Bildschirm, Umhang, Zahn, Mütze, Zelle, Hafen, Tonhöhe, Drachen, Socke, Schalter, Skala, Rücken, Schwingen, China, Wirbelsäule, Stamm, Hufeisen, Lehrer, Truthahn} \\
            }
        }
    }
    & \\ \\

    \theutterance \stepcounter{utterance}  
    & & & \multicolumn{4}{p{0.6\linewidth}}{
        \cellcolor[rgb]{0.9,0.9,0.9}{
            \makecell[{{p{\linewidth}}}]{
                \texttt{\tiny{[GM$\langle$P2]}}
                \texttt{TIPP: Zelle, Rücken, Wirbelsäule, Zahn} \\
            }
        }
    }
     \\ \\

    \theutterance \stepcounter{utterance}  
    & & & \multicolumn{2}{p{0.3\linewidth}}{
        \cellcolor[rgb]{0.9,0.9,0.9}{
            \makecell[{{p{\linewidth}}}]{
                \texttt{\tiny{[GM$|$GM]}}
                ['zelle', 'rücken', 'wirbelsäule', 'zahn']
            }
        }
    }
    & & \\ \\

    \theutterance \stepcounter{utterance}  
    & & & \multicolumn{2}{p{0.3\linewidth}}{
        \cellcolor[rgb]{0.9,0.9,0.9}{
            \makecell[{{p{\linewidth}}}]{
                \texttt{\tiny{[GM$|$GM]}}
                {'word': 'zelle', 'assignment': 'team'}
            }
        }
    }
    & & \\ \\

    \theutterance \stepcounter{utterance}  
    & & & \multicolumn{2}{p{0.3\linewidth}}{
        \cellcolor[rgb]{0.9,0.9,0.9}{
            \makecell[{{p{\linewidth}}}]{
                \texttt{\tiny{[GM$|$GM]}}
                {'word': 'zelle', 'assignment': 'team'}
            }
        }
    }
    & & \\ \\

    \theutterance \stepcounter{utterance}  
    & & & \multicolumn{2}{p{0.3\linewidth}}{
        \cellcolor[rgb]{0.9,0.9,0.9}{
            \makecell[{{p{\linewidth}}}]{
                \texttt{\tiny{[GM$|$GM]}}
                {'word': 'rücken', 'assignment': 'team'}
            }
        }
    }
    & & \\ \\

    \theutterance \stepcounter{utterance}  
    & & & \multicolumn{2}{p{0.3\linewidth}}{
        \cellcolor[rgb]{0.9,0.9,0.9}{
            \makecell[{{p{\linewidth}}}]{
                \texttt{\tiny{[GM$|$GM]}}
                {'word': 'rücken', 'assignment': 'team'}
            }
        }
    }
    & & \\ \\

    \theutterance \stepcounter{utterance}  
    & & & \multicolumn{2}{p{0.3\linewidth}}{
        \cellcolor[rgb]{0.9,0.9,0.9}{
            \makecell[{{p{\linewidth}}}]{
                \texttt{\tiny{[GM$|$GM]}}
                {'word': 'wirbelsäule', 'assignment': 'team'}
            }
        }
    }
    & & \\ \\

    \theutterance \stepcounter{utterance}  
    & & & \multicolumn{2}{p{0.3\linewidth}}{
        \cellcolor[rgb]{0.9,0.9,0.9}{
            \makecell[{{p{\linewidth}}}]{
                \texttt{\tiny{[GM$|$GM]}}
                {'word': 'wirbelsäule', 'assignment': 'team'}
            }
        }
    }
    & & \\ \\

    \theutterance \stepcounter{utterance}  
    & & & \multicolumn{2}{p{0.3\linewidth}}{
        \cellcolor[rgb]{0.9,0.9,0.9}{
            \makecell[{{p{\linewidth}}}]{
                \texttt{\tiny{[GM$|$GM]}}
                {'word': 'zahn', 'assignment': 'innocent'}
            }
        }
    }
    & & \\ \\

    \theutterance \stepcounter{utterance}  
    & & & \multicolumn{2}{p{0.3\linewidth}}{
        \cellcolor[rgb]{0.9,0.9,0.9}{
            \makecell[{{p{\linewidth}}}]{
                \texttt{\tiny{[GM$|$GM]}}
                \texttt{zahn} \\
            }
        }
    }
    & & \\ \\

    \theutterance \stepcounter{utterance}  
    & & & \multicolumn{2}{p{0.3\linewidth}}{
        \cellcolor[rgb]{0.9,0.9,0.9}{
            \makecell[{{p{\linewidth}}}]{
                \texttt{\tiny{[GM$|$GM]}}
                {'hidden': {'team': ['Hai', 'Stamm', 'Krawatte', 'Bildschirm', 'Schwingen', 'Hufeisen'], 'innocent': ['Schalter', 'Tonhöhe', 'China', 'Lehrer', 'Tokio', 'Drachen'], 'opponent': ['Hafen', 'Socke', 'Mitarbeiter', 'Mütze', 'Haken', 'Truthahn', 'Umhang', 'Skala'], 'assassin': ['Schule']}, 'revealed': {'team': {'team': ['Zelle', 'Rücken', 'Wirbelsäule'], 'innocent': ['Zahn'], 'opponent': [], 'assassin': []}, 'opponent': {'team': [], 'innocent': [], 'opponent': [], 'assassin': []}}}
            }
        }
    }
    & & \\ \\

    \theutterance \stepcounter{utterance}  
    & & \multicolumn{4}{p{0.6\linewidth}}{
        \cellcolor[rgb]{0.9,0.9,0.9}{
            \makecell[{{p{\linewidth}}}]{
                \texttt{\tiny{[P1$\langle$GM]}}
                \texttt{Die Wörter zelle, rücken, wirbelsäule wurden korrekt erraten. Das Wort zahn wurde erraten, ist aber ein unschuldig Wort. Der Zug deines Teamkollegen endete dort.} \\
\texttt{Gib jetzt einen weiteren Hinweis, der sich auf einige deiner verbleibenden Teamwörter und eine Liste der dazugehörigen Zielwörter bezieht. Denk daran, deinen Hinweis mit 'HINWEIS: ' zu beginnen, eine neue Zeile zu setzen und deine kommagetrennte Liste der Zielwörter mit 'ZIELE: ' zu beginnen. Hinweis: Im Vergleich zu vorherigen Anfragen wurden einige Wörter aus den Listen entfernt.} \\
\\ 
\texttt{Deine verbleibenden Teamwörter sind: Hai, Stamm, Krawatte, Bildschirm, Schwingen, Hufeisen.} \\
\texttt{Verbleibende Wörter für deinen Gegner sind: Hafen, Socke, Mitarbeiter, Mütze, Haken, Truthahn, Umhang, Skala.} \\
\texttt{Verbleibende Ablenkungswörter sind: Schalter, Tonhöhe, China, Lehrer, Tokio, Drachen.} \\
\texttt{Verbleibende Killerwörter sind: Schule.} \\
            }
        }
    }
    & \\ \\

    \theutterance \stepcounter{utterance}  
    & \multicolumn{4}{p{0.6\linewidth}}{
        \cellcolor[rgb]{0.9,0.9,0.9}{
            \makecell[{{p{\linewidth}}}]{
                \texttt{\tiny{[P1$\rangle$GM]}}
                \texttt{HINWEIS: Halfter} \\
\texttt{ZIELE: Hai, Krawatte} \\
            }
        }
    }
    & & \\ \\

    \theutterance \stepcounter{utterance}  
    & & & \multicolumn{2}{p{0.3\linewidth}}{
        \cellcolor[rgb]{0.9,0.9,0.9}{
            \makecell[{{p{\linewidth}}}]{
                \texttt{\tiny{[GM$|$GM]}}
                \texttt{halfter} \\
            }
        }
    }
    & & \\ \\

    \theutterance \stepcounter{utterance}  
    & & & \multicolumn{2}{p{0.3\linewidth}}{
        \cellcolor[rgb]{0.9,0.9,0.9}{
            \makecell[{{p{\linewidth}}}]{
                \texttt{\tiny{[GM$|$GM]}}
                ['hai', 'krawatte']
            }
        }
    }
    & & \\ \\

    \theutterance \stepcounter{utterance}  
    & & & \multicolumn{2}{p{0.3\linewidth}}{
        \cellcolor[rgb]{0.9,0.9,0.9}{
            \makecell[{{p{\linewidth}}}]{
                \texttt{\tiny{[GM$|$GM]}}
                {'word': 'hai', 'assignment': 'team'}
            }
        }
    }
    & & \\ \\

    \theutterance \stepcounter{utterance}  
    & & & \multicolumn{2}{p{0.3\linewidth}}{
        \cellcolor[rgb]{0.9,0.9,0.9}{
            \makecell[{{p{\linewidth}}}]{
                \texttt{\tiny{[GM$|$GM]}}
                {'word': 'krawatte', 'assignment': 'team'}
            }
        }
    }
    & & \\ \\

    \theutterance \stepcounter{utterance}  
    & & \multicolumn{4}{p{0.6\linewidth}}{
        \cellcolor[rgb]{0.9,0.9,0.9}{
            \makecell[{{p{\linewidth}}}]{
                \texttt{\tiny{[GM$\rangle$P2]}}
                \texttt{Die Wörter zelle, rücken, wirbelsäule wurden korrekt erraten. Das Wort zahn wurde erraten, ist aber ein unschuldig Wort. Dein Zug endete dort.} \\
\texttt{Gib jetzt eine weitere durch Kommas getrennte Liste mit mindestens 1 und bis zu 2 Wörtern aus der folgenden Wortliste an, die sich am besten auf das Wort 'halfter' beziehen oder am engsten damit verbunden sind. Denk daran, deine Antwort mit TIPP:  zu beginnen. Hinweis: Einige Wörter wurden im Vergleich zu vorherigen Anfragen aus der Liste entfernt.} \\
\\ 
\texttt{Schule, Tokio, Krawatte, Mitarbeiter, Hai, Haken, Bildschirm, Umhang, Mütze, Hafen, Tonhöhe, Drachen, Socke, Schalter, Skala, Schwingen, China, Stamm, Hufeisen, Lehrer, Truthahn} \\
            }
        }
    }
    & \\ \\

    \theutterance \stepcounter{utterance}  
    & & & \multicolumn{4}{p{0.6\linewidth}}{
        \cellcolor[rgb]{0.9,0.9,0.9}{
            \makecell[{{p{\linewidth}}}]{
                \texttt{\tiny{[GM$\langle$P2]}}
                \texttt{TIPP: Hufeisen, Haken} \\
            }
        }
    }
     \\ \\

    \theutterance \stepcounter{utterance}  
    & & & \multicolumn{2}{p{0.3\linewidth}}{
        \cellcolor[rgb]{0.9,0.9,0.9}{
            \makecell[{{p{\linewidth}}}]{
                \texttt{\tiny{[GM$|$GM]}}
                ['hufeisen', 'haken']
            }
        }
    }
    & & \\ \\

    \theutterance \stepcounter{utterance}  
    & & & \multicolumn{2}{p{0.3\linewidth}}{
        \cellcolor[rgb]{0.9,0.9,0.9}{
            \makecell[{{p{\linewidth}}}]{
                \texttt{\tiny{[GM$|$GM]}}
                {'word': 'hufeisen', 'assignment': 'team'}
            }
        }
    }
    & & \\ \\

    \theutterance \stepcounter{utterance}  
    & & & \multicolumn{2}{p{0.3\linewidth}}{
        \cellcolor[rgb]{0.9,0.9,0.9}{
            \makecell[{{p{\linewidth}}}]{
                \texttt{\tiny{[GM$|$GM]}}
                {'word': 'haken', 'assignment': 'opponent'}
            }
        }
    }
    & & \\ \\

    \theutterance \stepcounter{utterance}  
    & & & \multicolumn{2}{p{0.3\linewidth}}{
        \cellcolor[rgb]{0.9,0.9,0.9}{
            \makecell[{{p{\linewidth}}}]{
                \texttt{\tiny{[GM$|$GM]}}
                \texttt{haken} \\
            }
        }
    }
    & & \\ \\

    \theutterance \stepcounter{utterance}  
    & & & \multicolumn{2}{p{0.3\linewidth}}{
        \cellcolor[rgb]{0.9,0.9,0.9}{
            \makecell[{{p{\linewidth}}}]{
                \texttt{\tiny{[GM$|$GM]}}
                {'hidden': {'team': ['Hai', 'Stamm', 'Krawatte', 'Bildschirm', 'Schwingen'], 'innocent': ['Schalter', 'Tonhöhe', 'China', 'Lehrer', 'Tokio', 'Drachen'], 'opponent': ['Hafen', 'Socke', 'Mitarbeiter', 'Mütze', 'Truthahn', 'Umhang', 'Skala'], 'assassin': ['Schule']}, 'revealed': {'team': {'team': ['Zelle', 'Rücken', 'Wirbelsäule', 'Hufeisen'], 'innocent': ['Zahn'], 'opponent': ['Haken'], 'assassin': []}, 'opponent': {'team': [], 'innocent': [], 'opponent': [], 'assassin': []}}}
            }
        }
    }
    & & \\ \\

    \theutterance \stepcounter{utterance}  
    & & \multicolumn{4}{p{0.6\linewidth}}{
        \cellcolor[rgb]{0.9,0.9,0.9}{
            \makecell[{{p{\linewidth}}}]{
                \texttt{\tiny{[P1$\langle$GM]}}
                \texttt{Das Wort hufeisen wurde korrekt erraten. Das Wort haken wurde erraten, ist aber ein Gegner Wort. Der Zug deines Teamkollegen endete dort.} \\
\texttt{Gib jetzt einen weiteren Hinweis, der sich auf einige deiner verbleibenden Teamwörter und eine Liste der dazugehörigen Zielwörter bezieht. Denk daran, deinen Hinweis mit 'HINWEIS: ' zu beginnen, eine neue Zeile zu setzen und deine kommagetrennte Liste der Zielwörter mit 'ZIELE: ' zu beginnen. Hinweis: Im Vergleich zu vorherigen Anfragen wurden einige Wörter aus den Listen entfernt.} \\
\\ 
\texttt{Deine verbleibenden Teamwörter sind: Hai, Stamm, Krawatte, Bildschirm, Schwingen.} \\
\texttt{Verbleibende Wörter für deinen Gegner sind: Hafen, Socke, Mitarbeiter, Mütze, Truthahn, Umhang, Skala.} \\
\texttt{Verbleibende Ablenkungswörter sind: Schalter, Tonhöhe, China, Lehrer, Tokio, Drachen.} \\
\texttt{Verbleibende Killerwörter sind: Schule.} \\
            }
        }
    }
    & \\ \\

    \theutterance \stepcounter{utterance}  
    & \multicolumn{4}{p{0.6\linewidth}}{
        \cellcolor[rgb]{0.9,0.9,0.9}{
            \makecell[{{p{\linewidth}}}]{
                \texttt{\tiny{[P1$\rangle$GM]}}
                \texttt{HINWEIS: Flach} \\
\texttt{ZIELE: Bildschirm, Hai} \\
            }
        }
    }
    & & \\ \\

    \theutterance \stepcounter{utterance}  
    & & & \multicolumn{2}{p{0.3\linewidth}}{
        \cellcolor[rgb]{0.9,0.9,0.9}{
            \makecell[{{p{\linewidth}}}]{
                \texttt{\tiny{[GM$|$GM]}}
                \texttt{flach} \\
            }
        }
    }
    & & \\ \\

    \theutterance \stepcounter{utterance}  
    & & & \multicolumn{2}{p{0.3\linewidth}}{
        \cellcolor[rgb]{0.9,0.9,0.9}{
            \makecell[{{p{\linewidth}}}]{
                \texttt{\tiny{[GM$|$GM]}}
                ['bildschirm', 'hai']
            }
        }
    }
    & & \\ \\

    \theutterance \stepcounter{utterance}  
    & & & \multicolumn{2}{p{0.3\linewidth}}{
        \cellcolor[rgb]{0.9,0.9,0.9}{
            \makecell[{{p{\linewidth}}}]{
                \texttt{\tiny{[GM$|$GM]}}
                {'word': 'bildschirm', 'assignment': 'team'}
            }
        }
    }
    & & \\ \\

    \theutterance \stepcounter{utterance}  
    & & & \multicolumn{2}{p{0.3\linewidth}}{
        \cellcolor[rgb]{0.9,0.9,0.9}{
            \makecell[{{p{\linewidth}}}]{
                \texttt{\tiny{[GM$|$GM]}}
                {'word': 'hai', 'assignment': 'team'}
            }
        }
    }
    & & \\ \\

    \theutterance \stepcounter{utterance}  
    & & \multicolumn{4}{p{0.6\linewidth}}{
        \cellcolor[rgb]{0.9,0.9,0.9}{
            \makecell[{{p{\linewidth}}}]{
                \texttt{\tiny{[GM$\rangle$P2]}}
                \texttt{Das Wort hufeisen wurde korrekt erraten. Das Wort haken wurde erraten, ist aber ein Gegner Wort. Dein Zug endete dort.} \\
\texttt{Gib jetzt eine weitere durch Kommas getrennte Liste mit mindestens 1 und bis zu 2 Wörtern aus der folgenden Wortliste an, die sich am besten auf das Wort 'flach' beziehen oder am engsten damit verbunden sind. Denk daran, deine Antwort mit TIPP:  zu beginnen. Hinweis: Einige Wörter wurden im Vergleich zu vorherigen Anfragen aus der Liste entfernt.} \\
\\ 
\texttt{Schule, Tokio, Krawatte, Mitarbeiter, Hai, Bildschirm, Umhang, Mütze, Hafen, Tonhöhe, Drachen, Socke, Schalter, Skala, Schwingen, China, Stamm, Lehrer, Truthahn} \\
            }
        }
    }
    & \\ \\

    \theutterance \stepcounter{utterance}  
    & & & \multicolumn{4}{p{0.6\linewidth}}{
        \cellcolor[rgb]{0.9,0.9,0.9}{
            \makecell[{{p{\linewidth}}}]{
                \texttt{\tiny{[GM$\langle$P2]}}
                \texttt{TIPP: Skala, Schalter} \\
            }
        }
    }
     \\ \\

    \theutterance \stepcounter{utterance}  
    & & & \multicolumn{2}{p{0.3\linewidth}}{
        \cellcolor[rgb]{0.9,0.9,0.9}{
            \makecell[{{p{\linewidth}}}]{
                \texttt{\tiny{[GM$|$GM]}}
                ['skala', 'schalter']
            }
        }
    }
    & & \\ \\

    \theutterance \stepcounter{utterance}  
    & & & \multicolumn{2}{p{0.3\linewidth}}{
        \cellcolor[rgb]{0.9,0.9,0.9}{
            \makecell[{{p{\linewidth}}}]{
                \texttt{\tiny{[GM$|$GM]}}
                {'word': 'skala', 'assignment': 'opponent'}
            }
        }
    }
    & & \\ \\

    \theutterance \stepcounter{utterance}  
    & & & \multicolumn{2}{p{0.3\linewidth}}{
        \cellcolor[rgb]{0.9,0.9,0.9}{
            \makecell[{{p{\linewidth}}}]{
                \texttt{\tiny{[GM$|$GM]}}
                \texttt{skala} \\
            }
        }
    }
    & & \\ \\

    \theutterance \stepcounter{utterance}  
    & & & \multicolumn{2}{p{0.3\linewidth}}{
        \cellcolor[rgb]{0.9,0.9,0.9}{
            \makecell[{{p{\linewidth}}}]{
                \texttt{\tiny{[GM$|$GM]}}
                {'hidden': {'team': ['Hai', 'Stamm', 'Krawatte', 'Bildschirm', 'Schwingen'], 'innocent': ['Schalter', 'Tonhöhe', 'China', 'Lehrer', 'Tokio', 'Drachen'], 'opponent': ['Hafen', 'Socke', 'Mitarbeiter', 'Mütze', 'Truthahn', 'Umhang'], 'assassin': ['Schule']}, 'revealed': {'team': {'team': ['Zelle', 'Rücken', 'Wirbelsäule', 'Hufeisen'], 'innocent': ['Zahn'], 'opponent': ['Haken', 'Skala'], 'assassin': []}, 'opponent': {'team': [], 'innocent': [], 'opponent': [], 'assassin': []}}}
            }
        }
    }
    & & \\ \\

    \theutterance \stepcounter{utterance}  
    & & \multicolumn{4}{p{0.6\linewidth}}{
        \cellcolor[rgb]{0.9,0.9,0.9}{
            \makecell[{{p{\linewidth}}}]{
                \texttt{\tiny{[P1$\langle$GM]}}
                \texttt{Das Wort skala wurde erraten, ist aber ein Gegner Wort. Der Zug deines Teamkollegen endete dort.} \\
\texttt{Gib jetzt einen weiteren Hinweis, der sich auf einige deiner verbleibenden Teamwörter und eine Liste der dazugehörigen Zielwörter bezieht. Denk daran, deinen Hinweis mit 'HINWEIS: ' zu beginnen, eine neue Zeile zu setzen und deine kommagetrennte Liste der Zielwörter mit 'ZIELE: ' zu beginnen. Hinweis: Im Vergleich zu vorherigen Anfragen wurden einige Wörter aus den Listen entfernt.} \\
\\ 
\texttt{Deine verbleibenden Teamwörter sind: Hai, Stamm, Krawatte, Bildschirm, Schwingen.} \\
\texttt{Verbleibende Wörter für deinen Gegner sind: Hafen, Socke, Mitarbeiter, Mütze, Truthahn, Umhang.} \\
\texttt{Verbleibende Ablenkungswörter sind: Schalter, Tonhöhe, China, Lehrer, Tokio, Drachen.} \\
\texttt{Verbleibende Killerwörter sind: Schule.} \\
            }
        }
    }
    & \\ \\

    \theutterance \stepcounter{utterance}  
    & \multicolumn{4}{p{0.6\linewidth}}{
        \cellcolor[rgb]{0.9,0.9,0.9}{
            \makecell[{{p{\linewidth}}}]{
                \texttt{\tiny{[P1$\rangle$GM]}}
                \texttt{HINWEIS: Flattern} \\
\texttt{ZIELE: Schwingen, Krawatte} \\
            }
        }
    }
    & & \\ \\

    \theutterance \stepcounter{utterance}  
    & & & \multicolumn{2}{p{0.3\linewidth}}{
        \cellcolor[rgb]{0.9,0.9,0.9}{
            \makecell[{{p{\linewidth}}}]{
                \texttt{\tiny{[GM$|$GM]}}
                \texttt{flattern} \\
            }
        }
    }
    & & \\ \\

    \theutterance \stepcounter{utterance}  
    & & & \multicolumn{2}{p{0.3\linewidth}}{
        \cellcolor[rgb]{0.9,0.9,0.9}{
            \makecell[{{p{\linewidth}}}]{
                \texttt{\tiny{[GM$|$GM]}}
                ['schwingen', 'krawatte']
            }
        }
    }
    & & \\ \\

    \theutterance \stepcounter{utterance}  
    & & & \multicolumn{2}{p{0.3\linewidth}}{
        \cellcolor[rgb]{0.9,0.9,0.9}{
            \makecell[{{p{\linewidth}}}]{
                \texttt{\tiny{[GM$|$GM]}}
                {'word': 'schwingen', 'assignment': 'team'}
            }
        }
    }
    & & \\ \\

    \theutterance \stepcounter{utterance}  
    & & & \multicolumn{2}{p{0.3\linewidth}}{
        \cellcolor[rgb]{0.9,0.9,0.9}{
            \makecell[{{p{\linewidth}}}]{
                \texttt{\tiny{[GM$|$GM]}}
                {'word': 'krawatte', 'assignment': 'team'}
            }
        }
    }
    & & \\ \\

    \theutterance \stepcounter{utterance}  
    & & \multicolumn{4}{p{0.6\linewidth}}{
        \cellcolor[rgb]{0.9,0.9,0.9}{
            \makecell[{{p{\linewidth}}}]{
                \texttt{\tiny{[GM$\rangle$P2]}}
                \texttt{Das Wort skala wurde erraten, ist aber ein Gegner Wort. Dein Zug endete dort.} \\
\texttt{Gib jetzt eine weitere durch Kommas getrennte Liste mit mindestens 1 und bis zu 2 Wörtern aus der folgenden Wortliste an, die sich am besten auf das Wort 'flattern' beziehen oder am engsten damit verbunden sind. Denk daran, deine Antwort mit TIPP:  zu beginnen. Hinweis: Einige Wörter wurden im Vergleich zu vorherigen Anfragen aus der Liste entfernt.} \\
\\ 
\texttt{Schule, Tokio, Krawatte, Mitarbeiter, Hai, Bildschirm, Umhang, Mütze, Hafen, Tonhöhe, Drachen, Socke, Schalter, Schwingen, China, Stamm, Lehrer, Truthahn} \\
            }
        }
    }
    & \\ \\

    \theutterance \stepcounter{utterance}  
    & & & \multicolumn{4}{p{0.6\linewidth}}{
        \cellcolor[rgb]{0.9,0.9,0.9}{
            \makecell[{{p{\linewidth}}}]{
                \texttt{\tiny{[GM$\langle$P2]}}
                \texttt{TIPP: Schwingen, Drachen} \\
            }
        }
    }
     \\ \\

    \theutterance \stepcounter{utterance}  
    & & & \multicolumn{2}{p{0.3\linewidth}}{
        \cellcolor[rgb]{0.9,0.9,0.9}{
            \makecell[{{p{\linewidth}}}]{
                \texttt{\tiny{[GM$|$GM]}}
                ['schwingen', 'drachen']
            }
        }
    }
    & & \\ \\

    \theutterance \stepcounter{utterance}  
    & & & \multicolumn{2}{p{0.3\linewidth}}{
        \cellcolor[rgb]{0.9,0.9,0.9}{
            \makecell[{{p{\linewidth}}}]{
                \texttt{\tiny{[GM$|$GM]}}
                {'word': 'schwingen', 'assignment': 'team'}
            }
        }
    }
    & & \\ \\

    \theutterance \stepcounter{utterance}  
    & & & \multicolumn{2}{p{0.3\linewidth}}{
        \cellcolor[rgb]{0.9,0.9,0.9}{
            \makecell[{{p{\linewidth}}}]{
                \texttt{\tiny{[GM$|$GM]}}
                {'word': 'schwingen', 'assignment': 'team'}
            }
        }
    }
    & & \\ \\

    \theutterance \stepcounter{utterance}  
    & & & \multicolumn{2}{p{0.3\linewidth}}{
        \cellcolor[rgb]{0.9,0.9,0.9}{
            \makecell[{{p{\linewidth}}}]{
                \texttt{\tiny{[GM$|$GM]}}
                {'word': 'drachen', 'assignment': 'innocent'}
            }
        }
    }
    & & \\ \\

    \theutterance \stepcounter{utterance}  
    & & & \multicolumn{2}{p{0.3\linewidth}}{
        \cellcolor[rgb]{0.9,0.9,0.9}{
            \makecell[{{p{\linewidth}}}]{
                \texttt{\tiny{[GM$|$GM]}}
                \texttt{drachen} \\
            }
        }
    }
    & & \\ \\

    \theutterance \stepcounter{utterance}  
    & & & \multicolumn{2}{p{0.3\linewidth}}{
        \cellcolor[rgb]{0.9,0.9,0.9}{
            \makecell[{{p{\linewidth}}}]{
                \texttt{\tiny{[GM$|$GM]}}
                {'hidden': {'team': ['Hai', 'Stamm', 'Krawatte', 'Bildschirm'], 'innocent': ['Schalter', 'Tonhöhe', 'China', 'Lehrer', 'Tokio'], 'opponent': ['Hafen', 'Socke', 'Mitarbeiter', 'Mütze', 'Truthahn', 'Umhang'], 'assassin': ['Schule']}, 'revealed': {'team': {'team': ['Zelle', 'Rücken', 'Wirbelsäule', 'Hufeisen', 'Schwingen'], 'innocent': ['Zahn', 'Drachen'], 'opponent': ['Haken', 'Skala'], 'assassin': []}, 'opponent': {'team': [], 'innocent': [], 'opponent': [], 'assassin': []}}}
            }
        }
    }
    & & \\ \\

    \theutterance \stepcounter{utterance}  
    & & \multicolumn{4}{p{0.6\linewidth}}{
        \cellcolor[rgb]{0.9,0.9,0.9}{
            \makecell[{{p{\linewidth}}}]{
                \texttt{\tiny{[P1$\langle$GM]}}
                \texttt{Das Wort schwingen wurde korrekt erraten. Das Wort drachen wurde erraten, ist aber ein unschuldig Wort. Der Zug deines Teamkollegen endete dort.} \\
\texttt{Gib jetzt einen weiteren Hinweis, der sich auf einige deiner verbleibenden Teamwörter und eine Liste der dazugehörigen Zielwörter bezieht. Denk daran, deinen Hinweis mit 'HINWEIS: ' zu beginnen, eine neue Zeile zu setzen und deine kommagetrennte Liste der Zielwörter mit 'ZIELE: ' zu beginnen. Hinweis: Im Vergleich zu vorherigen Anfragen wurden einige Wörter aus den Listen entfernt.} \\
\\ 
\texttt{Deine verbleibenden Teamwörter sind: Hai, Stamm, Krawatte, Bildschirm.} \\
\texttt{Verbleibende Wörter für deinen Gegner sind: Hafen, Socke, Mitarbeiter, Mütze, Truthahn, Umhang.} \\
\texttt{Verbleibende Ablenkungswörter sind: Schalter, Tonhöhe, China, Lehrer, Tokio.} \\
\texttt{Verbleibende Killerwörter sind: Schule.} \\
            }
        }
    }
    & \\ \\

    \theutterance \stepcounter{utterance}  
    & \multicolumn{4}{p{0.6\linewidth}}{
        \cellcolor[rgb]{0.9,0.9,0.9}{
            \makecell[{{p{\linewidth}}}]{
                \texttt{\tiny{[P1$\rangle$GM]}}
                \texttt{HINWEIS: Kragen} \\
\texttt{ZIELE: Krawatte, Hai} \\
            }
        }
    }
    & & \\ \\

    \theutterance \stepcounter{utterance}  
    & & & \multicolumn{2}{p{0.3\linewidth}}{
        \cellcolor[rgb]{0.9,0.9,0.9}{
            \makecell[{{p{\linewidth}}}]{
                \texttt{\tiny{[GM$|$GM]}}
                \texttt{kragen} \\
            }
        }
    }
    & & \\ \\

    \theutterance \stepcounter{utterance}  
    & & & \multicolumn{2}{p{0.3\linewidth}}{
        \cellcolor[rgb]{0.9,0.9,0.9}{
            \makecell[{{p{\linewidth}}}]{
                \texttt{\tiny{[GM$|$GM]}}
                ['krawatte', 'hai']
            }
        }
    }
    & & \\ \\

    \theutterance \stepcounter{utterance}  
    & & & \multicolumn{2}{p{0.3\linewidth}}{
        \cellcolor[rgb]{0.9,0.9,0.9}{
            \makecell[{{p{\linewidth}}}]{
                \texttt{\tiny{[GM$|$GM]}}
                {'word': 'krawatte', 'assignment': 'team'}
            }
        }
    }
    & & \\ \\

    \theutterance \stepcounter{utterance}  
    & & & \multicolumn{2}{p{0.3\linewidth}}{
        \cellcolor[rgb]{0.9,0.9,0.9}{
            \makecell[{{p{\linewidth}}}]{
                \texttt{\tiny{[GM$|$GM]}}
                {'word': 'hai', 'assignment': 'team'}
            }
        }
    }
    & & \\ \\

    \theutterance \stepcounter{utterance}  
    & & \multicolumn{4}{p{0.6\linewidth}}{
        \cellcolor[rgb]{0.9,0.9,0.9}{
            \makecell[{{p{\linewidth}}}]{
                \texttt{\tiny{[GM$\rangle$P2]}}
                \texttt{Das Wort schwingen wurde korrekt erraten. Das Wort drachen wurde erraten, ist aber ein unschuldig Wort. Dein Zug endete dort.} \\
\texttt{Gib jetzt eine weitere durch Kommas getrennte Liste mit mindestens 1 und bis zu 2 Wörtern aus der folgenden Wortliste an, die sich am besten auf das Wort 'kragen' beziehen oder am engsten damit verbunden sind. Denk daran, deine Antwort mit TIPP:  zu beginnen. Hinweis: Einige Wörter wurden im Vergleich zu vorherigen Anfragen aus der Liste entfernt.} \\
\\ 
\texttt{Schule, Tokio, Krawatte, Mitarbeiter, Hai, Bildschirm, Umhang, Mütze, Hafen, Tonhöhe, Socke, Schalter, China, Stamm, Lehrer, Truthahn} \\
            }
        }
    }
    & \\ \\

    \theutterance \stepcounter{utterance}  
    & & & \multicolumn{4}{p{0.6\linewidth}}{
        \cellcolor[rgb]{0.9,0.9,0.9}{
            \makecell[{{p{\linewidth}}}]{
                \texttt{\tiny{[GM$\langle$P2]}}
                \texttt{TIPP: Umhang, Krawatte} \\
            }
        }
    }
     \\ \\

    \theutterance \stepcounter{utterance}  
    & & & \multicolumn{2}{p{0.3\linewidth}}{
        \cellcolor[rgb]{0.9,0.9,0.9}{
            \makecell[{{p{\linewidth}}}]{
                \texttt{\tiny{[GM$|$GM]}}
                ['umhang', 'krawatte']
            }
        }
    }
    & & \\ \\

    \theutterance \stepcounter{utterance}  
    & & & \multicolumn{2}{p{0.3\linewidth}}{
        \cellcolor[rgb]{0.9,0.9,0.9}{
            \makecell[{{p{\linewidth}}}]{
                \texttt{\tiny{[GM$|$GM]}}
                {'word': 'umhang', 'assignment': 'opponent'}
            }
        }
    }
    & & \\ \\

    \theutterance \stepcounter{utterance}  
    & & & \multicolumn{2}{p{0.3\linewidth}}{
        \cellcolor[rgb]{0.9,0.9,0.9}{
            \makecell[{{p{\linewidth}}}]{
                \texttt{\tiny{[GM$|$GM]}}
                \texttt{umhang} \\
            }
        }
    }
    & & \\ \\

    \theutterance \stepcounter{utterance}  
    & & & \multicolumn{2}{p{0.3\linewidth}}{
        \cellcolor[rgb]{0.9,0.9,0.9}{
            \makecell[{{p{\linewidth}}}]{
                \texttt{\tiny{[GM$|$GM]}}
                {'hidden': {'team': ['Hai', 'Stamm', 'Krawatte', 'Bildschirm'], 'innocent': ['Schalter', 'Tonhöhe', 'China', 'Lehrer', 'Tokio'], 'opponent': ['Hafen', 'Socke', 'Mitarbeiter', 'Mütze', 'Truthahn'], 'assassin': ['Schule']}, 'revealed': {'team': {'team': ['Zelle', 'Rücken', 'Wirbelsäule', 'Hufeisen', 'Schwingen'], 'innocent': ['Zahn', 'Drachen'], 'opponent': ['Haken', 'Skala', 'Umhang'], 'assassin': []}, 'opponent': {'team': [], 'innocent': [], 'opponent': [], 'assassin': []}}}
            }
        }
    }
    & & \\ \\

    \theutterance \stepcounter{utterance}  
    & & \multicolumn{4}{p{0.6\linewidth}}{
        \cellcolor[rgb]{0.9,0.9,0.9}{
            \makecell[{{p{\linewidth}}}]{
                \texttt{\tiny{[P1$\langle$GM]}}
                \texttt{Das Wort umhang wurde erraten, ist aber ein Gegner Wort. Der Zug deines Teamkollegen endete dort.} \\
\texttt{Gib jetzt einen weiteren Hinweis, der sich auf einige deiner verbleibenden Teamwörter und eine Liste der dazugehörigen Zielwörter bezieht. Denk daran, deinen Hinweis mit 'HINWEIS: ' zu beginnen, eine neue Zeile zu setzen und deine kommagetrennte Liste der Zielwörter mit 'ZIELE: ' zu beginnen. Hinweis: Im Vergleich zu vorherigen Anfragen wurden einige Wörter aus den Listen entfernt.} \\
\\ 
\texttt{Deine verbleibenden Teamwörter sind: Hai, Stamm, Krawatte, Bildschirm.} \\
\texttt{Verbleibende Wörter für deinen Gegner sind: Hafen, Socke, Mitarbeiter, Mütze, Truthahn.} \\
\texttt{Verbleibende Ablenkungswörter sind: Schalter, Tonhöhe, China, Lehrer, Tokio.} \\
\texttt{Verbleibende Killerwörter sind: Schule.} \\
            }
        }
    }
    & \\ \\

    \theutterance \stepcounter{utterance}  
    & \multicolumn{4}{p{0.6\linewidth}}{
        \cellcolor[rgb]{0.9,0.9,0.9}{
            \makecell[{{p{\linewidth}}}]{
                \texttt{\tiny{[P1$\rangle$GM]}}
                \texttt{HINWEIS: Monitor} \\
\texttt{ZIELE: Bildschirm} \\
            }
        }
    }
    & & \\ \\

    \theutterance \stepcounter{utterance}  
    & & & \multicolumn{2}{p{0.3\linewidth}}{
        \cellcolor[rgb]{0.9,0.9,0.9}{
            \makecell[{{p{\linewidth}}}]{
                \texttt{\tiny{[GM$|$GM]}}
                \texttt{monitor} \\
            }
        }
    }
    & & \\ \\

    \theutterance \stepcounter{utterance}  
    & & & \multicolumn{2}{p{0.3\linewidth}}{
        \cellcolor[rgb]{0.9,0.9,0.9}{
            \makecell[{{p{\linewidth}}}]{
                \texttt{\tiny{[GM$|$GM]}}
                ['bildschirm']
            }
        }
    }
    & & \\ \\

    \theutterance \stepcounter{utterance}  
    & & & \multicolumn{2}{p{0.3\linewidth}}{
        \cellcolor[rgb]{0.9,0.9,0.9}{
            \makecell[{{p{\linewidth}}}]{
                \texttt{\tiny{[GM$|$GM]}}
                {'word': 'bildschirm', 'assignment': 'team'}
            }
        }
    }
    & & \\ \\

    \theutterance \stepcounter{utterance}  
    & & \multicolumn{4}{p{0.6\linewidth}}{
        \cellcolor[rgb]{0.9,0.9,0.9}{
            \makecell[{{p{\linewidth}}}]{
                \texttt{\tiny{[GM$\rangle$P2]}}
                \texttt{Das Wort umhang wurde erraten, ist aber ein Gegner Wort. Dein Zug endete dort.} \\
\texttt{Gib jetzt eine weitere durch Kommas getrennte Liste mit mindestens 1 und bis zu 1 Wörtern aus der folgenden Wortliste an, die sich am besten auf das Wort 'monitor' beziehen oder am engsten damit verbunden sind. Denk daran, deine Antwort mit TIPP:  zu beginnen. Hinweis: Einige Wörter wurden im Vergleich zu vorherigen Anfragen aus der Liste entfernt.} \\
\\ 
\texttt{Schule, Tokio, Krawatte, Mitarbeiter, Hai, Bildschirm, Mütze, Hafen, Tonhöhe, Socke, Schalter, China, Stamm, Lehrer, Truthahn} \\
            }
        }
    }
    & \\ \\

    \theutterance \stepcounter{utterance}  
    & & & \multicolumn{4}{p{0.6\linewidth}}{
        \cellcolor[rgb]{0.9,0.9,0.9}{
            \makecell[{{p{\linewidth}}}]{
                \texttt{\tiny{[GM$\langle$P2]}}
                \texttt{TIPP: Bildschirm} \\
            }
        }
    }
     \\ \\

    \theutterance \stepcounter{utterance}  
    & & & \multicolumn{2}{p{0.3\linewidth}}{
        \cellcolor[rgb]{0.9,0.9,0.9}{
            \makecell[{{p{\linewidth}}}]{
                \texttt{\tiny{[GM$|$GM]}}
                ['bildschirm']
            }
        }
    }
    & & \\ \\

    \theutterance \stepcounter{utterance}  
    & & & \multicolumn{2}{p{0.3\linewidth}}{
        \cellcolor[rgb]{0.9,0.9,0.9}{
            \makecell[{{p{\linewidth}}}]{
                \texttt{\tiny{[GM$|$GM]}}
                {'word': 'bildschirm', 'assignment': 'team'}
            }
        }
    }
    & & \\ \\

    \theutterance \stepcounter{utterance}  
    & & & \multicolumn{2}{p{0.3\linewidth}}{
        \cellcolor[rgb]{0.9,0.9,0.9}{
            \makecell[{{p{\linewidth}}}]{
                \texttt{\tiny{[GM$|$GM]}}
                {'word': 'bildschirm', 'assignment': 'team'}
            }
        }
    }
    & & \\ \\

    \theutterance \stepcounter{utterance}  
    & & & \multicolumn{2}{p{0.3\linewidth}}{
        \cellcolor[rgb]{0.9,0.9,0.9}{
            \makecell[{{p{\linewidth}}}]{
                \texttt{\tiny{[GM$|$GM]}}
                {'hidden': {'team': ['Hai', 'Stamm', 'Krawatte'], 'innocent': ['Schalter', 'Tonhöhe', 'China', 'Lehrer', 'Tokio'], 'opponent': ['Hafen', 'Socke', 'Mitarbeiter', 'Mütze', 'Truthahn'], 'assassin': ['Schule']}, 'revealed': {'team': {'team': ['Zelle', 'Rücken', 'Wirbelsäule', 'Hufeisen', 'Schwingen', 'Bildschirm'], 'innocent': ['Zahn', 'Drachen'], 'opponent': ['Haken', 'Skala', 'Umhang'], 'assassin': []}, 'opponent': {'team': [], 'innocent': [], 'opponent': [], 'assassin': []}}}
            }
        }
    }
    & & \\ \\

    \theutterance \stepcounter{utterance}  
    & & \multicolumn{4}{p{0.6\linewidth}}{
        \cellcolor[rgb]{0.9,0.9,0.9}{
            \makecell[{{p{\linewidth}}}]{
                \texttt{\tiny{[P1$\langle$GM]}}
                \texttt{Das Wort bildschirm wurde korrekt erraten. Der Zug deines Teamkollegen endete dort.} \\
\texttt{Gib jetzt einen weiteren Hinweis, der sich auf einige deiner verbleibenden Teamwörter und eine Liste der dazugehörigen Zielwörter bezieht. Denk daran, deinen Hinweis mit 'HINWEIS: ' zu beginnen, eine neue Zeile zu setzen und deine kommagetrennte Liste der Zielwörter mit 'ZIELE: ' zu beginnen. Hinweis: Im Vergleich zu vorherigen Anfragen wurden einige Wörter aus den Listen entfernt.} \\
\\ 
\texttt{Deine verbleibenden Teamwörter sind: Hai, Stamm, Krawatte.} \\
\texttt{Verbleibende Wörter für deinen Gegner sind: Hafen, Socke, Mitarbeiter, Mütze, Truthahn.} \\
\texttt{Verbleibende Ablenkungswörter sind: Schalter, Tonhöhe, China, Lehrer, Tokio.} \\
\texttt{Verbleibende Killerwörter sind: Schule.} \\
            }
        }
    }
    & \\ \\

    \theutterance \stepcounter{utterance}  
    & \multicolumn{4}{p{0.6\linewidth}}{
        \cellcolor[rgb]{0.9,0.9,0.9}{
            \makecell[{{p{\linewidth}}}]{
                \texttt{\tiny{[P1$\rangle$GM]}}
                \texttt{HINWEIS: Knoten} \\
\texttt{ZIELE: Krawatte, Stamm} \\
            }
        }
    }
    & & \\ \\

    \theutterance \stepcounter{utterance}  
    & & & \multicolumn{2}{p{0.3\linewidth}}{
        \cellcolor[rgb]{0.9,0.9,0.9}{
            \makecell[{{p{\linewidth}}}]{
                \texttt{\tiny{[GM$|$GM]}}
                \texttt{knoten} \\
            }
        }
    }
    & & \\ \\

    \theutterance \stepcounter{utterance}  
    & & & \multicolumn{2}{p{0.3\linewidth}}{
        \cellcolor[rgb]{0.9,0.9,0.9}{
            \makecell[{{p{\linewidth}}}]{
                \texttt{\tiny{[GM$|$GM]}}
                ['krawatte', 'stamm']
            }
        }
    }
    & & \\ \\

    \theutterance \stepcounter{utterance}  
    & & & \multicolumn{2}{p{0.3\linewidth}}{
        \cellcolor[rgb]{0.9,0.9,0.9}{
            \makecell[{{p{\linewidth}}}]{
                \texttt{\tiny{[GM$|$GM]}}
                {'word': 'krawatte', 'assignment': 'team'}
            }
        }
    }
    & & \\ \\

    \theutterance \stepcounter{utterance}  
    & & & \multicolumn{2}{p{0.3\linewidth}}{
        \cellcolor[rgb]{0.9,0.9,0.9}{
            \makecell[{{p{\linewidth}}}]{
                \texttt{\tiny{[GM$|$GM]}}
                {'word': 'stamm', 'assignment': 'team'}
            }
        }
    }
    & & \\ \\

    \theutterance \stepcounter{utterance}  
    & & \multicolumn{4}{p{0.6\linewidth}}{
        \cellcolor[rgb]{0.9,0.9,0.9}{
            \makecell[{{p{\linewidth}}}]{
                \texttt{\tiny{[GM$\rangle$P2]}}
                \texttt{Das Wort bildschirm wurde korrekt erraten. Dein Zug endete dort.} \\
\texttt{Gib jetzt eine weitere durch Kommas getrennte Liste mit mindestens 1 und bis zu 2 Wörtern aus der folgenden Wortliste an, die sich am besten auf das Wort 'knoten' beziehen oder am engsten damit verbunden sind. Denk daran, deine Antwort mit TIPP:  zu beginnen. Hinweis: Einige Wörter wurden im Vergleich zu vorherigen Anfragen aus der Liste entfernt.} \\
\\ 
\texttt{Schule, Tokio, Krawatte, Mitarbeiter, Hai, Mütze, Hafen, Tonhöhe, Socke, Schalter, China, Stamm, Lehrer, Truthahn} \\
            }
        }
    }
    & \\ \\

    \theutterance \stepcounter{utterance}  
    & & & \multicolumn{4}{p{0.6\linewidth}}{
        \cellcolor[rgb]{0.9,0.9,0.9}{
            \makecell[{{p{\linewidth}}}]{
                \texttt{\tiny{[GM$\langle$P2]}}
                \texttt{TIPP: Krawatte, Schalter} \\
            }
        }
    }
     \\ \\

    \theutterance \stepcounter{utterance}  
    & & & \multicolumn{2}{p{0.3\linewidth}}{
        \cellcolor[rgb]{0.9,0.9,0.9}{
            \makecell[{{p{\linewidth}}}]{
                \texttt{\tiny{[GM$|$GM]}}
                ['krawatte', 'schalter']
            }
        }
    }
    & & \\ \\

    \theutterance \stepcounter{utterance}  
    & & & \multicolumn{2}{p{0.3\linewidth}}{
        \cellcolor[rgb]{0.9,0.9,0.9}{
            \makecell[{{p{\linewidth}}}]{
                \texttt{\tiny{[GM$|$GM]}}
                {'word': 'krawatte', 'assignment': 'team'}
            }
        }
    }
    & & \\ \\

    \theutterance \stepcounter{utterance}  
    & & & \multicolumn{2}{p{0.3\linewidth}}{
        \cellcolor[rgb]{0.9,0.9,0.9}{
            \makecell[{{p{\linewidth}}}]{
                \texttt{\tiny{[GM$|$GM]}}
                {'word': 'krawatte', 'assignment': 'team'}
            }
        }
    }
    & & \\ \\

    \theutterance \stepcounter{utterance}  
    & & & \multicolumn{2}{p{0.3\linewidth}}{
        \cellcolor[rgb]{0.9,0.9,0.9}{
            \makecell[{{p{\linewidth}}}]{
                \texttt{\tiny{[GM$|$GM]}}
                {'word': 'schalter', 'assignment': 'innocent'}
            }
        }
    }
    & & \\ \\

    \theutterance \stepcounter{utterance}  
    & & & \multicolumn{2}{p{0.3\linewidth}}{
        \cellcolor[rgb]{0.9,0.9,0.9}{
            \makecell[{{p{\linewidth}}}]{
                \texttt{\tiny{[GM$|$GM]}}
                \texttt{schalter} \\
            }
        }
    }
    & & \\ \\

    \theutterance \stepcounter{utterance}  
    & & & \multicolumn{2}{p{0.3\linewidth}}{
        \cellcolor[rgb]{0.9,0.9,0.9}{
            \makecell[{{p{\linewidth}}}]{
                \texttt{\tiny{[GM$|$GM]}}
                {'hidden': {'team': ['Hai', 'Stamm'], 'innocent': ['Tonhöhe', 'China', 'Lehrer', 'Tokio'], 'opponent': ['Hafen', 'Socke', 'Mitarbeiter', 'Mütze', 'Truthahn'], 'assassin': ['Schule']}, 'revealed': {'team': {'team': ['Zelle', 'Rücken', 'Wirbelsäule', 'Hufeisen', 'Schwingen', 'Bildschirm', 'Krawatte'], 'innocent': ['Zahn', 'Drachen', 'Schalter'], 'opponent': ['Haken', 'Skala', 'Umhang'], 'assassin': []}, 'opponent': {'team': [], 'innocent': [], 'opponent': [], 'assassin': []}}}
            }
        }
    }
    & & \\ \\

    \theutterance \stepcounter{utterance}  
    & & \multicolumn{4}{p{0.6\linewidth}}{
        \cellcolor[rgb]{0.9,0.9,0.9}{
            \makecell[{{p{\linewidth}}}]{
                \texttt{\tiny{[P1$\langle$GM]}}
                \texttt{Das Wort krawatte wurde korrekt erraten. Das Wort schalter wurde erraten, ist aber ein unschuldig Wort. Der Zug deines Teamkollegen endete dort.} \\
\texttt{Gib jetzt einen weiteren Hinweis, der sich auf einige deiner verbleibenden Teamwörter und eine Liste der dazugehörigen Zielwörter bezieht. Denk daran, deinen Hinweis mit 'HINWEIS: ' zu beginnen, eine neue Zeile zu setzen und deine kommagetrennte Liste der Zielwörter mit 'ZIELE: ' zu beginnen. Hinweis: Im Vergleich zu vorherigen Anfragen wurden einige Wörter aus den Listen entfernt.} \\
\\ 
\texttt{Deine verbleibenden Teamwörter sind: Hai, Stamm.} \\
\texttt{Verbleibende Wörter für deinen Gegner sind: Hafen, Socke, Mitarbeiter, Mütze, Truthahn.} \\
\texttt{Verbleibende Ablenkungswörter sind: Tonhöhe, China, Lehrer, Tokio.} \\
\texttt{Verbleibende Killerwörter sind: Schule.} \\
            }
        }
    }
    & \\ \\

    \theutterance \stepcounter{utterance}  
    & \multicolumn{4}{p{0.6\linewidth}}{
        \cellcolor[rgb]{0.9,0.9,0.9}{
            \makecell[{{p{\linewidth}}}]{
                \texttt{\tiny{[P1$\rangle$GM]}}
                \texttt{HINWEIS: Blutsverwandtschaft} \\
\texttt{ZIELE: Stamm, Hai} \\
            }
        }
    }
    & & \\ \\

    \theutterance \stepcounter{utterance}  
    & & & \multicolumn{2}{p{0.3\linewidth}}{
        \cellcolor[rgb]{0.9,0.9,0.9}{
            \makecell[{{p{\linewidth}}}]{
                \texttt{\tiny{[GM$|$GM]}}
                \texttt{blutsverwandtschaft} \\
            }
        }
    }
    & & \\ \\

    \theutterance \stepcounter{utterance}  
    & & & \multicolumn{2}{p{0.3\linewidth}}{
        \cellcolor[rgb]{0.9,0.9,0.9}{
            \makecell[{{p{\linewidth}}}]{
                \texttt{\tiny{[GM$|$GM]}}
                ['stamm', 'hai']
            }
        }
    }
    & & \\ \\

    \theutterance \stepcounter{utterance}  
    & & & \multicolumn{2}{p{0.3\linewidth}}{
        \cellcolor[rgb]{0.9,0.9,0.9}{
            \makecell[{{p{\linewidth}}}]{
                \texttt{\tiny{[GM$|$GM]}}
                {'word': 'stamm', 'assignment': 'team'}
            }
        }
    }
    & & \\ \\

    \theutterance \stepcounter{utterance}  
    & & & \multicolumn{2}{p{0.3\linewidth}}{
        \cellcolor[rgb]{0.9,0.9,0.9}{
            \makecell[{{p{\linewidth}}}]{
                \texttt{\tiny{[GM$|$GM]}}
                {'word': 'hai', 'assignment': 'team'}
            }
        }
    }
    & & \\ \\

    \theutterance \stepcounter{utterance}  
    & & \multicolumn{4}{p{0.6\linewidth}}{
        \cellcolor[rgb]{0.9,0.9,0.9}{
            \makecell[{{p{\linewidth}}}]{
                \texttt{\tiny{[GM$\rangle$P2]}}
                \texttt{Das Wort krawatte wurde korrekt erraten. Das Wort schalter wurde erraten, ist aber ein unschuldig Wort. Dein Zug endete dort.} \\
\texttt{Gib jetzt eine weitere durch Kommas getrennte Liste mit mindestens 1 und bis zu 2 Wörtern aus der folgenden Wortliste an, die sich am besten auf das Wort 'blutsverwandtschaft' beziehen oder am engsten damit verbunden sind. Denk daran, deine Antwort mit TIPP:  zu beginnen. Hinweis: Einige Wörter wurden im Vergleich zu vorherigen Anfragen aus der Liste entfernt.} \\
\\ 
\texttt{Schule, Tokio, Mitarbeiter, Hai, Mütze, Hafen, Tonhöhe, Socke, China, Stamm, Lehrer, Truthahn} \\
            }
        }
    }
    & \\ \\

    \theutterance \stepcounter{utterance}  
    & & & \multicolumn{4}{p{0.6\linewidth}}{
        \cellcolor[rgb]{0.9,0.9,0.9}{
            \makecell[{{p{\linewidth}}}]{
                \texttt{\tiny{[GM$\langle$P2]}}
                \texttt{TIPP: Stamm, China} \\
            }
        }
    }
     \\ \\

    \theutterance \stepcounter{utterance}  
    & & & \multicolumn{2}{p{0.3\linewidth}}{
        \cellcolor[rgb]{0.9,0.9,0.9}{
            \makecell[{{p{\linewidth}}}]{
                \texttt{\tiny{[GM$|$GM]}}
                ['stamm', 'china']
            }
        }
    }
    & & \\ \\

    \theutterance \stepcounter{utterance}  
    & & & \multicolumn{2}{p{0.3\linewidth}}{
        \cellcolor[rgb]{0.9,0.9,0.9}{
            \makecell[{{p{\linewidth}}}]{
                \texttt{\tiny{[GM$|$GM]}}
                {'word': 'stamm', 'assignment': 'team'}
            }
        }
    }
    & & \\ \\

    \theutterance \stepcounter{utterance}  
    & & & \multicolumn{2}{p{0.3\linewidth}}{
        \cellcolor[rgb]{0.9,0.9,0.9}{
            \makecell[{{p{\linewidth}}}]{
                \texttt{\tiny{[GM$|$GM]}}
                {'word': 'stamm', 'assignment': 'team'}
            }
        }
    }
    & & \\ \\

    \theutterance \stepcounter{utterance}  
    & & & \multicolumn{2}{p{0.3\linewidth}}{
        \cellcolor[rgb]{0.9,0.9,0.9}{
            \makecell[{{p{\linewidth}}}]{
                \texttt{\tiny{[GM$|$GM]}}
                {'word': 'china', 'assignment': 'innocent'}
            }
        }
    }
    & & \\ \\

    \theutterance \stepcounter{utterance}  
    & & & \multicolumn{2}{p{0.3\linewidth}}{
        \cellcolor[rgb]{0.9,0.9,0.9}{
            \makecell[{{p{\linewidth}}}]{
                \texttt{\tiny{[GM$|$GM]}}
                \texttt{china} \\
            }
        }
    }
    & & \\ \\

    \theutterance \stepcounter{utterance}  
    & & & \multicolumn{2}{p{0.3\linewidth}}{
        \cellcolor[rgb]{0.9,0.9,0.9}{
            \makecell[{{p{\linewidth}}}]{
                \texttt{\tiny{[GM$|$GM]}}
                {'hidden': {'team': ['Hai'], 'innocent': ['Tonhöhe', 'Lehrer', 'Tokio'], 'opponent': ['Hafen', 'Socke', 'Mitarbeiter', 'Mütze', 'Truthahn'], 'assassin': ['Schule']}, 'revealed': {'team': {'team': ['Zelle', 'Rücken', 'Wirbelsäule', 'Hufeisen', 'Schwingen', 'Bildschirm', 'Krawatte', 'Stamm'], 'innocent': ['Zahn', 'Drachen', 'Schalter', 'China'], 'opponent': ['Haken', 'Skala', 'Umhang'], 'assassin': []}, 'opponent': {'team': [], 'innocent': [], 'opponent': [], 'assassin': []}}}
            }
        }
    }
    & & \\ \\

    \theutterance \stepcounter{utterance}  
    & & \multicolumn{4}{p{0.6\linewidth}}{
        \cellcolor[rgb]{0.9,0.9,0.9}{
            \makecell[{{p{\linewidth}}}]{
                \texttt{\tiny{[P1$\langle$GM]}}
                \texttt{Das Wort stamm wurde korrekt erraten. Das Wort china wurde erraten, ist aber ein unschuldig Wort. Der Zug deines Teamkollegen endete dort.} \\
\texttt{Gib jetzt einen weiteren Hinweis, der sich auf einige deiner verbleibenden Teamwörter und eine Liste der dazugehörigen Zielwörter bezieht. Denk daran, deinen Hinweis mit 'HINWEIS: ' zu beginnen, eine neue Zeile zu setzen und deine kommagetrennte Liste der Zielwörter mit 'ZIELE: ' zu beginnen. Hinweis: Im Vergleich zu vorherigen Anfragen wurden einige Wörter aus den Listen entfernt.} \\
\\ 
\texttt{Deine verbleibenden Teamwörter sind: Hai.} \\
\texttt{Verbleibende Wörter für deinen Gegner sind: Hafen, Socke, Mitarbeiter, Mütze, Truthahn.} \\
\texttt{Verbleibende Ablenkungswörter sind: Tonhöhe, Lehrer, Tokio.} \\
\texttt{Verbleibende Killerwörter sind: Schule.} \\
            }
        }
    }
    & \\ \\

    \theutterance \stepcounter{utterance}  
    & \multicolumn{4}{p{0.6\linewidth}}{
        \cellcolor[rgb]{0.9,0.9,0.9}{
            \makecell[{{p{\linewidth}}}]{
                \texttt{\tiny{[P1$\rangle$GM]}}
                \texttt{HINWEIS: Meeresräuber} \\
\texttt{ZIELE: Hai} \\
            }
        }
    }
    & & \\ \\

    \theutterance \stepcounter{utterance}  
    & & & \multicolumn{2}{p{0.3\linewidth}}{
        \cellcolor[rgb]{0.9,0.9,0.9}{
            \makecell[{{p{\linewidth}}}]{
                \texttt{\tiny{[GM$|$GM]}}
                \texttt{meeresräuber} \\
            }
        }
    }
    & & \\ \\

    \theutterance \stepcounter{utterance}  
    & & & \multicolumn{2}{p{0.3\linewidth}}{
        \cellcolor[rgb]{0.9,0.9,0.9}{
            \makecell[{{p{\linewidth}}}]{
                \texttt{\tiny{[GM$|$GM]}}
                ['hai']
            }
        }
    }
    & & \\ \\

    \theutterance \stepcounter{utterance}  
    & & & \multicolumn{2}{p{0.3\linewidth}}{
        \cellcolor[rgb]{0.9,0.9,0.9}{
            \makecell[{{p{\linewidth}}}]{
                \texttt{\tiny{[GM$|$GM]}}
                {'word': 'hai', 'assignment': 'team'}
            }
        }
    }
    & & \\ \\

    \theutterance \stepcounter{utterance}  
    & & \multicolumn{4}{p{0.6\linewidth}}{
        \cellcolor[rgb]{0.9,0.9,0.9}{
            \makecell[{{p{\linewidth}}}]{
                \texttt{\tiny{[GM$\rangle$P2]}}
                \texttt{Das Wort stamm wurde korrekt erraten. Das Wort china wurde erraten, ist aber ein unschuldig Wort. Dein Zug endete dort.} \\
\texttt{Gib jetzt eine weitere durch Kommas getrennte Liste mit mindestens 1 und bis zu 1 Wörtern aus der folgenden Wortliste an, die sich am besten auf das Wort 'meeresräuber' beziehen oder am engsten damit verbunden sind. Denk daran, deine Antwort mit TIPP:  zu beginnen. Hinweis: Einige Wörter wurden im Vergleich zu vorherigen Anfragen aus der Liste entfernt.} \\
\\ 
\texttt{Schule, Tokio, Mitarbeiter, Hai, Mütze, Hafen, Tonhöhe, Socke, Lehrer, Truthahn} \\
            }
        }
    }
    & \\ \\

    \theutterance \stepcounter{utterance}  
    & & & \multicolumn{4}{p{0.6\linewidth}}{
        \cellcolor[rgb]{0.9,0.9,0.9}{
            \makecell[{{p{\linewidth}}}]{
                \texttt{\tiny{[GM$\langle$P2]}}
                \texttt{TIPP: Hafen} \\
            }
        }
    }
     \\ \\

    \theutterance \stepcounter{utterance}  
    & & & \multicolumn{2}{p{0.3\linewidth}}{
        \cellcolor[rgb]{0.9,0.9,0.9}{
            \makecell[{{p{\linewidth}}}]{
                \texttt{\tiny{[GM$|$GM]}}
                ['hafen']
            }
        }
    }
    & & \\ \\

    \theutterance \stepcounter{utterance}  
    & & & \multicolumn{2}{p{0.3\linewidth}}{
        \cellcolor[rgb]{0.9,0.9,0.9}{
            \makecell[{{p{\linewidth}}}]{
                \texttt{\tiny{[GM$|$GM]}}
                {'word': 'hafen', 'assignment': 'opponent'}
            }
        }
    }
    & & \\ \\

    \theutterance \stepcounter{utterance}  
    & & & \multicolumn{2}{p{0.3\linewidth}}{
        \cellcolor[rgb]{0.9,0.9,0.9}{
            \makecell[{{p{\linewidth}}}]{
                \texttt{\tiny{[GM$|$GM]}}
                \texttt{hafen} \\
            }
        }
    }
    & & \\ \\

    \theutterance \stepcounter{utterance}  
    & & & \multicolumn{2}{p{0.3\linewidth}}{
        \cellcolor[rgb]{0.9,0.9,0.9}{
            \makecell[{{p{\linewidth}}}]{
                \texttt{\tiny{[GM$|$GM]}}
                {'hidden': {'team': ['Hai'], 'innocent': ['Tonhöhe', 'Lehrer', 'Tokio'], 'opponent': ['Socke', 'Mitarbeiter', 'Mütze', 'Truthahn'], 'assassin': ['Schule']}, 'revealed': {'team': {'team': ['Zelle', 'Rücken', 'Wirbelsäule', 'Hufeisen', 'Schwingen', 'Bildschirm', 'Krawatte', 'Stamm'], 'innocent': ['Zahn', 'Drachen', 'Schalter', 'China'], 'opponent': ['Haken', 'Skala', 'Umhang', 'Hafen'], 'assassin': []}, 'opponent': {'team': [], 'innocent': [], 'opponent': [], 'assassin': []}}}
            }
        }
    }
    & & \\ \\

    \theutterance \stepcounter{utterance}  
    & & \multicolumn{4}{p{0.6\linewidth}}{
        \cellcolor[rgb]{0.9,0.9,0.9}{
            \makecell[{{p{\linewidth}}}]{
                \texttt{\tiny{[P1$\langle$GM]}}
                \texttt{Das Wort hafen wurde erraten, ist aber ein Gegner Wort. Der Zug deines Teamkollegen endete dort.} \\
\texttt{Gib jetzt einen weiteren Hinweis, der sich auf einige deiner verbleibenden Teamwörter und eine Liste der dazugehörigen Zielwörter bezieht. Denk daran, deinen Hinweis mit 'HINWEIS: ' zu beginnen, eine neue Zeile zu setzen und deine kommagetrennte Liste der Zielwörter mit 'ZIELE: ' zu beginnen. Hinweis: Im Vergleich zu vorherigen Anfragen wurden einige Wörter aus den Listen entfernt.} \\
\\ 
\texttt{Deine verbleibenden Teamwörter sind: Hai.} \\
\texttt{Verbleibende Wörter für deinen Gegner sind: Socke, Mitarbeiter, Mütze, Truthahn.} \\
\texttt{Verbleibende Ablenkungswörter sind: Tonhöhe, Lehrer, Tokio.} \\
\texttt{Verbleibende Killerwörter sind: Schule.} \\
            }
        }
    }
    & \\ \\

    \theutterance \stepcounter{utterance}  
    & \multicolumn{4}{p{0.6\linewidth}}{
        \cellcolor[rgb]{0.9,0.9,0.9}{
            \makecell[{{p{\linewidth}}}]{
                \texttt{\tiny{[P1$\rangle$GM]}}
                \texttt{HINWEIS: Fisch} \\
\texttt{ZIELE: Hai} \\
            }
        }
    }
    & & \\ \\

    \theutterance \stepcounter{utterance}  
    & & & \multicolumn{2}{p{0.3\linewidth}}{
        \cellcolor[rgb]{0.9,0.9,0.9}{
            \makecell[{{p{\linewidth}}}]{
                \texttt{\tiny{[GM$|$GM]}}
                \texttt{fisch} \\
            }
        }
    }
    & & \\ \\

    \theutterance \stepcounter{utterance}  
    & & & \multicolumn{2}{p{0.3\linewidth}}{
        \cellcolor[rgb]{0.9,0.9,0.9}{
            \makecell[{{p{\linewidth}}}]{
                \texttt{\tiny{[GM$|$GM]}}
                ['hai']
            }
        }
    }
    & & \\ \\

    \theutterance \stepcounter{utterance}  
    & & & \multicolumn{2}{p{0.3\linewidth}}{
        \cellcolor[rgb]{0.9,0.9,0.9}{
            \makecell[{{p{\linewidth}}}]{
                \texttt{\tiny{[GM$|$GM]}}
                {'word': 'hai', 'assignment': 'team'}
            }
        }
    }
    & & \\ \\

    \theutterance \stepcounter{utterance}  
    & & \multicolumn{4}{p{0.6\linewidth}}{
        \cellcolor[rgb]{0.9,0.9,0.9}{
            \makecell[{{p{\linewidth}}}]{
                \texttt{\tiny{[GM$\rangle$P2]}}
                \texttt{Das Wort hafen wurde erraten, ist aber ein Gegner Wort. Dein Zug endete dort.} \\
\texttt{Gib jetzt eine weitere durch Kommas getrennte Liste mit mindestens 1 und bis zu 1 Wörtern aus der folgenden Wortliste an, die sich am besten auf das Wort 'fisch' beziehen oder am engsten damit verbunden sind. Denk daran, deine Antwort mit TIPP:  zu beginnen. Hinweis: Einige Wörter wurden im Vergleich zu vorherigen Anfragen aus der Liste entfernt.} \\
\\ 
\texttt{Schule, Tokio, Mitarbeiter, Hai, Mütze, Tonhöhe, Socke, Lehrer, Truthahn} \\
            }
        }
    }
    & \\ \\

    \theutterance \stepcounter{utterance}  
    & & & \multicolumn{4}{p{0.6\linewidth}}{
        \cellcolor[rgb]{0.9,0.9,0.9}{
            \makecell[{{p{\linewidth}}}]{
                \texttt{\tiny{[GM$\langle$P2]}}
                \texttt{TIPP: Hai} \\
            }
        }
    }
     \\ \\

    \theutterance \stepcounter{utterance}  
    & & & \multicolumn{2}{p{0.3\linewidth}}{
        \cellcolor[rgb]{0.9,0.9,0.9}{
            \makecell[{{p{\linewidth}}}]{
                \texttt{\tiny{[GM$|$GM]}}
                ['hai']
            }
        }
    }
    & & \\ \\

    \theutterance \stepcounter{utterance}  
    & & & \multicolumn{2}{p{0.3\linewidth}}{
        \cellcolor[rgb]{0.9,0.9,0.9}{
            \makecell[{{p{\linewidth}}}]{
                \texttt{\tiny{[GM$|$GM]}}
                {'word': 'hai', 'assignment': 'team'}
            }
        }
    }
    & & \\ \\

    \theutterance \stepcounter{utterance}  
    & & & \multicolumn{2}{p{0.3\linewidth}}{
        \cellcolor[rgb]{0.9,0.9,0.9}{
            \makecell[{{p{\linewidth}}}]{
                \texttt{\tiny{[GM$|$GM]}}
                {'word': 'hai', 'assignment': 'team'}
            }
        }
    }
    & & \\ \\

    \theutterance \stepcounter{utterance}  
    & & & \multicolumn{2}{p{0.3\linewidth}}{
        \cellcolor[rgb]{0.9,0.9,0.9}{
            \makecell[{{p{\linewidth}}}]{
                \texttt{\tiny{[GM$|$GM]}}
                \texttt{team has won} \\
            }
        }
    }
    & & \\ \\

\end{supertabular}
}

\end{document}
